\subsection*{Justificacion Fila 6:}

Primero se renombra la fila A de la tabla R, por el nombre C.

\begin{center}
$\rho_{C \leftarrow A}(R) = R'=$ \quad
\begin{tabular}{|c|c|}
  \hline
  \textbf{C} & \textbf{B} \\ \hline
  1 & x \\
  2 & y \\
  2 & z \\
  3 & x \\
  9 & a \\ \hline
\end{tabular}
\end{center}

Posteriormente se aplica Join Natural sobre la tabla anterior y la tabla S. Dado que las tablas comparten las columnas C y B, entonces con el join natural buscamos las tuplas en donde los valores de estas columnas coinciden. Notemos que solo coinciden para $C = 2$ y $B = y$). Así, tenemos la siguiente tabla.
\begin{center}
$R' \bowtie S =$ \quad
\begin{tabular}{|c|c|c|}
  \hline
 
  \textbf{C} & \textbf{B} & \textbf{D} \\ \hline
  2 & y & 1 \\ \hline
\end{tabular}
\end{center}
